\begin{tikzpicture}[font=\small]
  % ---------- baseband DT spectrum, 2pi-periodic ----------
  \begin{scope}
    \draw[ax] (-6.4,0) -- (6.8,0); \node[note, anchor=north west] at (6.6,-0.05){$\Omega$};
    \draw[ax] (0,-0.2) -- (0,1.8);
    \node[note, anchor=south east] at (-0.12,1.4){$X(e^{j\Omega})$};
    \foreach \c in {-6.0,0,6.0}{
      \draw[curve] (\c-1.1,0) -- (\c,1.35) -- (\c+1.1,0);
    }
    \draw[helper] (3.0,-0.2) -- (3.0,1.6);
    \draw[helper] (-3.0,-0.2) -- (-3.0,1.6);
    \node[note, anchor=north] at (3.0,-0.1){$\pi$};
    \node[note, anchor=north] at (6.0,-0.1){$2\pi$};
    \node[note, anchor=north] at (-3.0,-0.1){$-\pi$};
    \node[note, anchor=north, align=center] at (0,-0.8){(a) 基带序列的谱，本身就以 $2\pi$ 为周期};
  \end{scope}

  % ---------- after modulation by cos(Omega_c n) ----------
  \begin{scope}[yshift=-3.4cm]
    \draw[ax] (-6.4,0) -- (6.8,0); \node[note, anchor=north west] at (6.6,-0.05){$\Omega$};
    \draw[ax] (0,-0.2) -- (0,1.8);
    \node[note, anchor=south east] at (-0.12,1.4){$Y(e^{j\Omega})$};
    \foreach \c in {-6.0,0,6.0}{
      \foreach \d in {-1.7,1.7}{
        \draw[curve3] (\c+\d-1.1,0) -- (\c+\d,0.68) -- (\c+\d+1.1,0);
      }
    }
    \draw[helper] (3.0,-0.2) -- (3.0,1.6);
    \draw[helper] (-3.0,-0.2) -- (-3.0,1.6);
    \node[note, anchor=north] at (3.0,-0.1){$\pi$};
    \node[note, anchor=north] at (1.7,-0.1){$\Omega_c$};
    \node[note, anchor=north] at (-3.0,-0.1){$-\pi$};
    \node[note, anchor=north, align=center] at (0,-0.8)
      {(b) 乘 $\cos(\Omega_cn)$ 后同样搬移 $\pm\Omega_c$，幅度减半，周期性仍在};
  \end{scope}

  \node[note, anchor=north, align=center] at (0,-5.5)
    {与连续时间调制\textbf{形式完全相同}，但多了一条硬约束：
     所有频率都住在长度 $2\pi$ 的区间里。\\
     所以 $\Omega_c$ 不能随便调大 —— 搬得太远，副本就会从 $\pi$ 那一头绕回来撞上自己。\\
     无混叠条件：$\Omega_c+\Omega_M<\pi$（单边带 / 复调制可以放宽这条）};
\end{tikzpicture}
